\documentclass[11pt]{article}

\usepackage{amsmath,amssymb,amsthm}
\usepackage[a4paper,margin=1in]{geometry}
\usepackage{hyperref}

% Remove paragraph indentation
\setlength{\parindent}{0pt}
\setlength{\parskip}{6pt}

\begin{document}

\title{A Concise Derivation of ISOMAP}
\author{Zhenwei Tang}
\date{}
\maketitle

% ─────────────────────────────────────────────
\section*{Goal}
% ─────────────────────────────────────────────

Given $n$ high-dimensional data points
\[
  x_i \in \mathbb{R}^p, \qquad i = 1, \dots, n,
\]
ISOMAP aims to find a low-dimensional embedding
\[
  z_i \in \mathbb{R}^k, \qquad k \ll p,
\]
such that the \emph{geodesic} (manifold) distances between points are preserved:
\[
  \|z_i - z_j\|_2 \approx d_{\mathcal{M}}(x_i, x_j),
\]
where $d_{\mathcal{M}}(x_i, x_j)$ denotes the geodesic distance along the underlying
manifold $\mathcal{M}$.

The key insight is that Euclidean distances in the ambient space $\mathbb{R}^p$ can
be misleading when the data lie on a curved manifold; geodesic distances faithfully
reflect the intrinsic geometry. ISOMAP approximates geodesic distances via shortest
paths on a neighbourhood graph, then recovers coordinates $z_i$ by applying
Classical Multidimensional Scaling (CMDS).

% ─────────────────────────────────────────────
\section*{Step 1: Construct the Neighbourhood Graph}
% ─────────────────────────────────────────────

\textbf{Why:} Geodesic distances cannot be computed directly; we approximate them by
"hopping" along nearby data points.

For each point $x_i$, connect it to its $K$ nearest neighbours (or to all points
within a ball of radius $\varepsilon$) using Euclidean distance. Assign edge weight
\[
  w_{ij} = \|x_i - x_j\|_2
\]
to each edge $(i,j)$ in the graph $G$. Let
\[
  W \in \mathbb{R}^{n \times n}, \quad W_{ij} =
  \begin{cases} \|x_i - x_j\|_2 & \text{if } (i,j) \in G, \\ \infty & \text{otherwise.} \end{cases}
\]

% ─────────────────────────────────────────────
\section*{Step 2: Compute All-Pairs Shortest Paths}
% ─────────────────────────────────────────────

\textbf{Why:} The geodesic distance between two points is the length of the shortest
path along the manifold. On the graph, this corresponds to the shortest weighted path.

\textbf{Using Dijkstra's algorithm} (or Floyd–Warshall) on $G$, compute the
geodesic-distance matrix
\[
  D_G \in \mathbb{R}^{n \times n}, \quad (D_G)_{ij} = d_G(x_i, x_j),
\]
where $d_G(x_i, x_j)$ is the length of the shortest path between $x_i$ and $x_j$ in
$G$. As $n \to \infty$ and $K$ (or $\varepsilon$) grows appropriately,
$d_G(x_i, x_j) \to d_{\mathcal{M}}(x_i, x_j)$.

% ─────────────────────────────────────────────
\section*{Step 3: Convert Distances to Inner Products (Double Centering)}
% ─────────────────────────────────────────────

\textbf{Why:} Classical MDS (applied in Step 4) operates on an inner-product matrix,
not a distance matrix. We use double centering to recover inner products from squared
distances.

Let $\delta_{ij} = (D_G)_{ij}$ denote the estimated geodesic distance. Define the
squared-distance matrix
\[
  \Delta \in \mathbb{R}^{n \times n}, \quad \Delta_{ij} = \delta_{ij}^2.
\]

\textbf{Applying the double-centering operator} (see Appendix~\ref{app:doublecenter}):

Define the centering matrix
\[
  H = I_n - \frac{1}{n}\mathbf{1}\mathbf{1}^T \in \mathbb{R}^{n \times n},
\]
where $I_n$ is the $n \times n$ identity matrix and $\mathbf{1} \in \mathbb{R}^n$ is
the all-ones vector.

Compute the inner-product (Gram) matrix:
\[
  \boxed{B = -\frac{1}{2} H \Delta H.}
\]

\textbf{Why this works:} If the embedding coordinates are $z_i \in \mathbb{R}^k$,
then $\|z_i - z_j\|^2 = \|z_i\|^2 - 2z_i^T z_j + \|z_j\|^2$. After double
centering (which removes the row/column means and thus eliminates the $\|z_i\|^2$
and $\|z_j\|^2$ terms), one recovers $B_{ij} = z_i^T z_j$, i.e.\ $B = ZZ^T$ where
$Z = [z_1|\cdots|z_n]^T$.

% ─────────────────────────────────────────────
\section*{Step 4: Eigendecomposition of $B$ (Classical MDS)}
% ─────────────────────────────────────────────

\textbf{Why:} We want to find coordinates $Z \in \mathbb{R}^{n \times k}$ such that
$ZZ^T = B$. This is a matrix factorization problem solved by eigendecomposition.

\textbf{Applying eigendecomposition} (see Appendix~\ref{app:eigen}) to the symmetric
matrix $B$:
\[
  B = V \Lambda V^T,
\]
where
\begin{align*}
  V &= [v_1 \mid v_2 \mid \cdots \mid v_n] \in \mathbb{R}^{n \times n}, \quad
      \text{orthonormal eigenvectors},\\
  \Lambda &= \operatorname{diag}(\lambda_1, \lambda_2, \dots, \lambda_n), \quad
      \lambda_1 \geq \lambda_2 \geq \cdots \geq \lambda_n.
\end{align*}

\textbf{Why $B$ is positive semi-definite:} If the geodesic distances are exactly
Euclidean in some space, then $B = ZZ^T$ is a Gram matrix and hence PSD (see
Appendix~\ref{app:psd}). In practice, $B$ may have small negative eigenvalues due to
approximation error; these are discarded.

Retain the top $k$ non-negative eigenvalues $\lambda_1 \geq \cdots \geq \lambda_k > 0$
and their eigenvectors $v_1, \dots, v_k$. Define
\[
  V_k = [v_1 \mid \cdots \mid v_k] \in \mathbb{R}^{n \times k}, \qquad
  \Lambda_k = \operatorname{diag}(\lambda_1, \dots, \lambda_k) \in \mathbb{R}^{k \times k}.
\]

% ─────────────────────────────────────────────
\section*{Step 5: Recover the Low-Dimensional Embedding}
% ─────────────────────────────────────────────

\textbf{Why:} We need $Z$ such that $ZZ^T \approx B$. Using the truncated
eigendecomposition $B \approx V_k \Lambda_k V_k^T$, we set
\[
  Z = V_k \Lambda_k^{1/2},
\]
where $\Lambda_k^{1/2} = \operatorname{diag}(\sqrt{\lambda_1}, \dots, \sqrt{\lambda_k})$.

\textbf{Verification:}
\[
  ZZ^T = V_k \Lambda_k^{1/2} \left(V_k \Lambda_k^{1/2}\right)^T
       = V_k \Lambda_k^{1/2} \Lambda_k^{1/2} V_k^T
       = V_k \Lambda_k V_k^T \approx B. \quad \checkmark
\]

The $i$-th row of $Z$ gives the embedding coordinate:
\[
  z_i = \Lambda_k^{1/2} V_k^T e_i \in \mathbb{R}^k,
\]
where $e_i$ is the $i$-th standard basis vector (i.e.\ the $i$-th row of $Z$).

% ─────────────────────────────────────────────
\section*{Conclusion}
% ─────────────────────────────────────────────

ISOMAP recovers the intrinsic low-dimensional geometry of data on a nonlinear
manifold through three main stages:

\begin{enumerate}
  \item \textbf{Graph construction:} approximate geodesic distances via shortest
        paths on the $K$-NN graph, obtaining $D_G$.
  \item \textbf{Double centering:} convert squared geodesic distances to inner
        products via $B = -\tfrac{1}{2}H\Delta H$.
  \item \textbf{Classical MDS:} recover embedding coordinates
        $Z = V_k\Lambda_k^{1/2}$ from the top-$k$ eigenpairs of $B$.
\end{enumerate}

The final embedding satisfies
\[
  \|z_i - z_j\|_2 \approx d_{\mathcal{M}}(x_i, x_j),
\]
i.e.\ Euclidean distances in the low-dimensional space faithfully approximate
geodesic distances on the original manifold.

\begin{center}
\fbox{\parbox{0.85\textwidth}{\centering
ISOMAP embeds high-dimensional manifold data into $\mathbb{R}^k$ by replacing
Euclidean distances with graph-approximated geodesic distances, then applying
Classical MDS to recover coordinates that preserve intrinsic geometry.
}}
\end{center}

% ─────────────────────────────────────────────
\appendix
\section{Double Centering}
\label{app:doublecenter}
% ─────────────────────────────────────────────

Let $z_1, \dots, z_n \in \mathbb{R}^k$ be centred embeddings, i.e.\
$\sum_i z_i = 0$. Define $\bar{z} = 0$. Then
\[
  \delta_{ij}^2 = \|z_i - z_j\|^2 = \|z_i\|^2 - 2z_i^T z_j + \|z_j\|^2.
\]

Let $a_i = \|z_i\|^2$. Then
\begin{align*}
  \frac{1}{n}\sum_j \delta_{ij}^2 &= a_i + \frac{1}{n}\sum_j a_j, \\
  \frac{1}{n}\sum_i \delta_{ij}^2 &= a_j + \frac{1}{n}\sum_i a_i, \\
  \frac{1}{n^2}\sum_{i,j} \delta_{ij}^2 &= \frac{2}{n}\sum_i a_i.
\end{align*}

The double-centering formula
\[
  B_{ij} = -\frac{1}{2}\left(\delta_{ij}^2
    - \frac{1}{n}\sum_{j'}\delta_{ij'}^2
    - \frac{1}{n}\sum_{i'}\delta_{i'j}^2
    + \frac{1}{n^2}\sum_{i',j'}\delta_{i'j'}^2\right)
\]
simplifies to
\[
  B_{ij} = -\frac{1}{2}\bigl(-2z_i^T z_j\bigr) = z_i^T z_j,
\]
i.e.\ $B = ZZ^T$. In matrix form, $B = -\tfrac{1}{2}H\Delta H$ where
$H = I - \tfrac{1}{n}\mathbf{1}\mathbf{1}^T$ is the \emph{centering matrix}
satisfying $H^2 = H$ and $H\mathbf{1} = 0$.

% ─────────────────────────────────────────────
\section{Eigendecomposition of a Symmetric Matrix}
\label{app:eigen}
% ─────────────────────────────────────────────

\textbf{Spectral Theorem:} For any real symmetric matrix $A = A^T \in \mathbb{R}^{n
\times n}$, there exists an orthogonal matrix $V \in \mathbb{R}^{n \times n}$
(i.e.\ $V^TV = I$) and a diagonal matrix $\Lambda =
\operatorname{diag}(\lambda_1,\dots,\lambda_n)$ with $\lambda_1 \geq \cdots \geq
\lambda_n$ such that
\[
  A = V\Lambda V^T.
\]
The columns $v_i$ of $V$ are the \emph{eigenvectors} of $A$, and $\lambda_i$ are
the corresponding \emph{eigenvalues}, satisfying $Av_i = \lambda_i v_i$.

% ─────────────────────────────────────────────
\section{Positive Semi-Definite (PSD) Matrix}
\label{app:psd}
% ─────────────────────────────────────────────

A symmetric matrix $A \in \mathbb{R}^{n\times n}$ is \emph{positive semi-definite}
(PSD), written $A \succeq 0$, if
\[
  u^T A u \geq 0 \quad \forall\, u \in \mathbb{R}^n.
\]

\textbf{Equivalently:} all eigenvalues of $A$ are non-negative, i.e.\ $\lambda_i
\geq 0$ for all $i$.

\textbf{Gram matrices are PSD:} For any $Z \in \mathbb{R}^{n \times k}$, the matrix
$B = ZZ^T$ satisfies
\[
  u^T B u = u^T ZZ^T u = \|Z^Tu\|^2 \geq 0,
\]
so $B \succeq 0$. This guarantees that, when geodesic distances are exactly
Euclidean-embeddable, $B$ has no negative eigenvalues and the ISOMAP embedding is
exact.

% ─────────────────────────────────────────────
\section{Dijkstra's Algorithm (Shortest Path)}
\label{app:dijkstra}
% ─────────────────────────────────────────────

Given a weighted graph $G = (V, E, w)$ with non-negative weights $w_{ij} \geq 0$,
\textbf{Dijkstra's algorithm} computes the shortest-path distance from a source node
$s$ to all other nodes in $O(|E| + |V|\log|V|)$ time. Applied $n$ times (once per
source), it yields the full all-pairs shortest-path matrix $D_G$ in
$O(n(|E| + n\log n))$ time. Alternatively, \textbf{Floyd–Warshall} computes all
pairs in $O(n^3)$ time, which is simpler to implement but less efficient for sparse
graphs.

\end{document}
